% !PRINTED 544
\weberpage{544}
Es existieren 63 negative Diskriminanten von der Eigenschaft, daß in jedem Geschlecht nur eine Klasse enthalten ist; daß es nicht mehr gibt, selbst daß die Zahl nur eine endliche ist, kann freilich bis jetzt nur durch Induktion geschlossen, nicht bewiesen werden. Die Mehrzahl dieser Diskriminanten, die bereits in \S\,139 zusammengestellt sind, ist $\equiv 1, 2\ (\mathrm{mod}\ 4)$ und ohne quadratischen Teiler oder $\equiv 8\ (\mathrm{mod}\ 16)$.

Für die ersteren lassen sich die \gesperrt{Klasseninvarianten} nach der Formel (10) vollständig berechnen, und eine ähnliche Formel, auf deren Bildung wir hier nicht eingehen wollen, führt auch für die durch 8 teilbaren Diskriminanten zum Ziel.

Für die vereinzelten Diskriminanten dieser Art, die hierher nicht passen, lassen sich die Klasseninvarianten $f(\omega)$ nach einer der anderen Methoden berechnen.\footnote{Vgl.\ des Verfassers Abhandlung: \glqq Zur komplexen Multiplikation elliptischer Funktionen\grqq{}. Mathematische Annalen, Bd.\,23. Ich mache hier auf einen Fehler in der \gesperrt{Gauss}schen Tafel aufmerksam, den ich bei Gelegenheit dieser Rechnungen gefunden habe: \gesperrt{Gauss'} Werke, Bd.\,II, S.\,475 muß die positive Determinante 136 die Bezeichnung IV, 2, nicht IV, 1 haben.}

Um an einem einfachen, leicht zu übersehenden Beispiele die Rechnung durchzuführen, wählen wir $m=105=3.\,5.\,7$. Wenn wir die Werte von $\delta$, die $\equiv 1\ (\mathrm{mod}\ 8)$ sind, weglassen, da diese in der Summe der Formeln (10) keinen Beitrag geben, so haben wir:
\[
\begin{aligned}
\delta &= -3, & 5, & & 21, & & -35, \\
\delta' &= 140, & -84, & & -20, & & 12
\end{aligned}
\]
zu setzen und erhalten, da $g=8$ ist:
% [sic] The source prints radicands 84 and 20 in rows 2 and 3; the positive discriminants 5 and 21 are required by (11) and are used explicitly on p.545.
\[
\begin{aligned}
16\log\frac{f(\sqrt{-105})}{\sqrt[4]{2}} &= K(-3)\,K(140)\log\frac{T+U\sqrt{140}}{2}, \\
&\quad + K(5)\,K(-84)\log\frac{T+U\sqrt{84}}{2}, \\
&\quad + K(21)\,K(-20)\log\frac{T+U\sqrt{20}}{2}, \\
&\quad + K(-35)\,K(12)\log\frac{T+U\sqrt{12}}{2}.
\end{aligned}
\]
